\section{Conclusion and Discussion}
\label{Section: Conclusion}
+\red{This paper developed DualHash, a stochastic primal--dual method for deep hashing with W-type regularization. The analysis first characterizes the convex envelope $h^{**}$ and its contact set, then gives an exact stationarity correspondence for KKT points of the conjugate formulation and a near-stationarity certificate for the original objective $J_0$. The certificate separates an optimization term controlled by the KKT residual from a structural term controlled by $h(\mathcal{D}(\bx))$; for the $\ell_1$ penalty, the latter is $\lambda Q_1(\bx)$. Under the stated assumptions, DualHash-StoM and DualHash-StoRM attain oracle complexities of $\mathcal{O}(\varepsilon^{-4})$ and $\mathcal{O}(\varepsilon^{-3})$, respectively, for $\varepsilon$-KKT points. The experiments report mAP, precision, and quantization error for CIFAR-10, NUS-WIDE, and ImageNet-100 across code lengths, with the observed rankings varying across datasets and bit lengths.}

\red{The same conjugate construction may also be useful in other composite problems with nonsmooth W-type penalties, but establishing such extensions requires separate analysis.}

% \begin{remark}\label{rmk:penalty-surrogate}
% For any fixed $\gamma > 0$, the penalty problem \eqref{Prop: block-wise regularized problem} 
% is a surrogate of the constrained formulation and is not equivalent to the 
% original problem \eqref{Eq: the primal deep hashing problem}. All convergence 
% results in this paper (Theorems~\ref{Theorem: the oracle complexity in SGDM} 
% and~\ref{Theorem: the oracle complexity in STORM}) are established for the 
% Lagrangian of~\eqref{Prop: block-wise regularized problem}. 
% A formal characterization of the gap between stationary points of 
% \eqref{Prop: block-wise regularized problem} and those of the original 
% composite problem as $\gamma \to \infty$ is an interesting direction 
% for future work.
% \end{remark}
